% ---------- Слайд 14: Overleaf vs Локальная компиляция ----------
\begin{frame}[label=tools]{Среды компиляции: Overleaf и локальная установка}
  \centering
  \renewcommand{\arraystretch}{1.3}
  \begin{tabular}{@{}p{0.45\textwidth} p{0.45\textwidth}@{}}
    \toprule
    \textbf{Overleaf (облачная среда)} & \textbf{Локальная установка} \\
    \midrule
    Отсутствие необходимости установки ПО & Максимальная скорость компиляции \\
    Доступ с любого устройства & Полный контроль над окружением \\
    Встроенные средства совместной работы & Независимость от сети Интернет \\
    Ограничения бесплатного тарифа & Требуется настройка дистрибутива \\
    \bottomrule
  \end{tabular}

  \vspace{1em}
  \begin{block}{Гибридный подход}
    Локальная работа с использованием системы контроля версий (Git)
    и синхронизация с облачной средой для совместной работы или
    доступа с различных устройств.
  \end{block}
  \note{Для первичного освоения рекомендуется использовать Overleaf
  ввиду минимального порога входа. Для подготовки диссертации
  и крупных проектов предпочтительна локальная установка
  (например, TeX Live в связке с редактором VS Code).}
\end{frame}

% ---------- Слайд 15: Типичные ошибки ----------
\begin{frame}[label=mistakes]{Типичные ошибки при подготовке слайдов}
  \begin{itemize}
    \item \textbf{Избыточный объем текста.} Рекомендуется формулировать один тезис на слайд.
    \item \textbf{Мелкий кегль шрифта.} Текст, формулы и код должны быть различимы с последнего ряда аудитории.
    \item \textbf{Неоправданное применение анимации.} Анимация используется исключительно для акцентирования логики изложения.
    \item \textbf{Использование растровой графики.} Экспорт диаграмм в PNG/JPEG приводит к потере качества при масштабировании. Необходимо использовать векторные форматы (PDF, EPS).
  \end{itemize}

  \vspace{0.5em}
  \begin{alertblock}{Практическое замечание}
    Перед выступлением рекомендуется провести компиляцию с включёнными
    заметками (\cmd{setbeameroption\{show notes\}}) и провести репетицию.
  \end{alertblock}
  \note{Данный перечень ошибок носит универсальный характер
  и относится к подготовке научных презентаций в целом,
  вне зависимости от используемого программного обеспечения.}
\end{frame}